% =============================================================================
%  book/ccs.tex  ——  全书导言区（版式 / 字体 / 自定义命令与环境）
%  版式依据：O'Reilly《Programming Rust, 2nd Edition》英文原版 PDF 实测
%     页面 504 × 661.5bp（7in × 9.19in）；版心 360bp（左右边距各 72bp）；
%     正文 Minion Pro 10.5pt / 行距 13.2pt；章号 16.8bp、章名 25.2bp、
%     节 18.9bp、小节 15.8bp（均为 MyriadPro-SemiboldCond，黑色）；
%     代码 UbuntuMono 8.5bp，无框无底色；页脚为「页码 | 章名 / 节名」+ 细线。
%  文件结构沿用 CMake-Cookbook 模板：
%     main.tex → book/ccs.tex（导言区）→ book/index.tex（封面 + 目录 + 章节装配）
%     正文：book/content/<类别>/<slug>.tex，由 book/index.tex 的 \myChapter 加载
% =============================================================================

\documentclass[10.5bp,openright,twoside,UTF8]{book}
\interfootnotelinepenalty=10000

% -----------------------------------------------------------------------------
%  基础宏包
% -----------------------------------------------------------------------------
\usepackage{textcomp}
\usepackage[T1]{fontenc}
\usepackage{microtype}
\usepackage{amsmath}
\usepackage{amssymb}
\usepackage{ifthen}
\usepackage{pgffor}
\usepackage{xspace}
\usepackage{graphicx}
\usepackage{tikz}
\usepackage{tabularx}
\usepackage{xltabular}   % longtable + X 列：行数多的表要跨页（见 table() 注释）
\usepackage{booktabs}
\usepackage{multirow}
\usepackage{float}
\usepackage{enumitem}
\usepackage{multicol}
\usepackage{indentfirst}
\usepackage{subfiles}
\usepackage[table]{xcolor}
\usepackage{caption}
\usepackage[export]{adjustbox}

% -----------------------------------------------------------------------------
%  ★★★ 页面尺寸与页边距：想改版面，只动这一段 ★★★
%  单位一律用 bp（1/72 in），与 PDF 同单位，便于和原书实测值对齐。
%    纸张     504 × 661.5bp
%    版心     360bp 宽（left = right = 72bp）
%    版心高   528bp（40 行 × 13.2bp）；正文首行基线 y ≈ 63bp
%    ⚠ 章首页的版式由 \titlespacing*{\chapter} 控制，改 top 会一起移动。
% -----------------------------------------------------------------------------
\usepackage[
  paperwidth=504bp, paperheight=661.5bp,
  textwidth=360bp,
  top=53bp, bottom=78.5bp,
  headheight=14bp, headsep=10bp,
  footskip=36bp]{geometry}

% -----------------------------------------------------------------------------
%  ★★★ 字体：想换字体、改字号，只改这一段 ★★★
%  原书用 Minion Pro（正文）/ MyriadPro-SemiboldCond（标题）/ UbuntuMono（代码），
%  三者都是商用字体，本机以 Cambria / Arial Narrow / Consolas 代替：
%    · Cambria 的数字是等高数字，x 高度与 Minion 接近，替换后行宽变化最小；
%    · Arial Narrow 是压缩无衬线，最接近 MyriadPro-SemiboldCond 的字宽密度；
%    · Consolas 自带全套框线字符（代码里的目录树 ─├└│ 可以直接显示）。
% -----------------------------------------------------------------------------
\usepackage{fontspec}
\usepackage[fontset=none]{ctex}      % 只借 ctex 的中文标题名与 \zihao，字体手工指定
\usepackage[fontsize=10.5bp]{fontsize}
\linespread{1.0476}                 % 10.5bp × 1.2 × 1.0476 ≈ 13.2bp 行距（原书实测）

\setCJKmainfont[AutoFakeSlant=true]{Source Han Serif SC}                 % 中文正文
\setCJKsansfont[AutoFakeBold=true,AutoFakeSlant=true]{Microsoft YaHei}   % 中文标题/表头
\setCJKmonofont[AutoFakeBold=true,AutoFakeSlant=true]{Microsoft YaHei}   % 代码里的中文
\setmainfont{Cambria}[BoldFont=*-Bold, ItalicFont=*-Italic, BoldItalicFont=*-BoldItalic]
\setsansfont{Arial Narrow}
\setmonofont{Consolas}[
  Path=C:/Windows/Fonts/, Extension=.ttf,
  UprightFont=consola, BoldFont=consolab,
  ItalicFont=consolai, BoldItalicFont=consolaz]
\newfontfamily\codehan{Microsoft YaHei}

% -----------------------------------------------------------------------------
%  版面与段落
% -----------------------------------------------------------------------------
\definecolor{ink}{HTML}{000000}          % 正文墨色（原书为纯黑）
\definecolor{linkred}{HTML}{990000}      % 交叉引用 / 链接（原书实测 #990000）
\definecolor{tabhead}{HTML}{000000}      % 表头黑底
\definecolor{figframe}{HTML}{BFBFBF}     % 插图细边框
\definecolor{rule}{HTML}{666666}         % 页脚横线

\setlength{\parindent}{2em}                      % 中文惯例：首行缩进两字
\setlength{\parskip}{0bp}
\raggedbottom
\emergencystretch=1em                            % 只给 1em 应急拉伸，避免中文字距被拉散
\AtBeginDocument{\color{ink}}

\setlist[itemize,enumerate,description]{nosep, topsep=5bp, partopsep=0bp,
  parsep=0bp, itemsep=3bp, leftmargin=2.2em}

% 脚注：原书 8bp，行距 10bp，编号为正文内上标
\renewcommand{\footnotesize}{\fontsize{8bp}{10bp}\selectfont}
% minus 必须留够：脚注页（正文末行 + 分隔线 + 脚注）一旦超出版心，
% TeX 只能报 Overfull \vbox 而不会自己挤，实测第 13 章有一页超 11.5bp。
% 允许压缩 12bp 后，这类页会被自动吸收，正常页不受影响。
\setlength{\skip\footins}{14bp plus 4bp minus 14bp}
\renewcommand{\footnoterule}{\kern-3bp\rule{120bp}{0.3bp}\kern 6bp}
\makeatletter
\renewcommand{\@makefnmark}{\hbox{\textsuperscript{\fontsize{7bp}{7bp}\selectfont\@thefnmark}}}
\makeatother

% -----------------------------------------------------------------------------
%  代码高亮
%  原书代码：无外框、无底色、无行号，8.5bp 等宽，左缘比正文再缩进 17bp。
%
%  ★ 两条路线，用下面的开关切换 ★
%    \mintedcodefalse（默认）—— 构建期用 Pygments 预着色：convert.py 直接把
%      \PYG{...} 序列写进 .tex，编译时只是 Verbatim 环境，不需要 -shell-escape，
%      也不需要 latexminted。全书 1400+ 代码块，单遍编译约 3 分钟。
%    \mintedcodetrue  —— 编译期用 minted 高亮（需 -shell-escape + Pygments + Python
%      在 PATH 里）。观感一致，但这台机器上单遍编译要 19 分钟，且 minted v3 在
%      Windows 下有缓存目录必须为绝对路径的坑（见下面的补丁，已保留备查）。
% -----------------------------------------------------------------------------
\newif\ifmintedcode
\mintedcodefalse
\newcommand{\codefont}{\fontsize{8.5bp}{10.2bp}\selectfont}

\ifmintedcode
\usepackage[cache=true]{minted}
\setminted{obeytabs, tabsize=4, breaklines=true, breakanywhere=true,
           fontsize=\codefont, autogobble=true}
% 【关键】minted v3 在 Windows 下缓存目录必须是绝对路径：
%   latexminted 只允许把文件写进 TEXMFOUTPUT / TEXMF_OUTPUT_DIRECTORY 指明的
%   绝对路径内，相对路径一律拒绝；而 minted 默认 \minted@cachedir = "_minted"
%   是相对路径 → 全部代码块写入失败，且报错信息极具误导性。
\makeatletter
\CatchFileDef{\mintedProjectRoot}{|"kpsewhich --var-value TEXMF_OUTPUT_DIRECTORY"}%
  {\minted@standardcatcodes\endlinechar=-1}%
\def\minted@stripcr#1\r{#1}%
\edef\mintedProjectRoot{\expandafter\minted@stripcr\mintedProjectRoot\r}%
\ifx\mintedProjectRoot\@empty
  \CatchFileDef{\mintedProjectRoot}{|"kpsewhich --var-value TEXMFOUTPUT"}%
    {\minted@standardcatcodes\endlinechar=-1}%
  \edef\mintedProjectRoot{\expandafter\minted@stripcr\mintedProjectRoot\r}%
\fi
\ifx\mintedProjectRoot\@empty
  \PackageWarning{minted-cn}{%
    Neither TEXMF_OUTPUT_DIRECTORY nor TEXMFOUTPUT is set: code blocks will
    FAIL to highlight. Build via build.bat.}%
\else
  \edef\minted@cachedir{\mintedProjectRoot}%
\fi
\makeatother
% 注册原书配色，避免每次去调 Python 生成样式
\makeatletter
\CatchFileDef{\minted@styledef@oreilly}{book/oreilly-pygments.tex}%
  {\minted@standardcatcodes\endlinechar=-1}%
\minted@detectconfig
\makeatother
\setminted{style=oreilly}

% 代码块：直接用 minted 环境（无框、无底色、无行号）
\newenvironment{codebox}[1][rust]
  {\VerbatimEnvironment
   \begin{minted}[breaklines=true, breakanywhere=true,
                  fontsize=\codefont, tabsize=4, obeytabs=true,
                  autogobble=true, xleftmargin=17bp]{#1}}
  {\end{minted}}

\else
% ---- 默认路线：构建期预着色 ----
%  fvextra 提供 Verbatim 的 breaklines / xleftmargin 等选项
\usepackage{fvextra}
% =============================================================================
%  O'Reilly《Programming Rust》代码配色（从英文原版 PDF 逐 token 采样）
%  关键字 #006699 粗体 · 类型 #007788 粗体 · 字符串 #cc3300 · 数字 #ff6600
%  注释 #336666 斜体 · 标点/运算符 #555555 · 名字空间 #000088 · 函数名 #cc00ff
%  由 book/ccs.tex 读入并注册为 minted 样式 "oreilly"（避免再次调用 Python 生成样式）
% =============================================================================
\makeatletter
\def\PYG@reset{\let\PYG@it=\relax \let\PYG@bf=\relax%
    \let\PYG@ul=\relax \let\PYG@tc=\relax%
    \let\PYG@bc=\relax \let\PYG@ff=\relax}
\def\PYG@tok#1{\csname PYG@tok@#1\endcsname}
\def\PYG@toks#1+{\ifx\relax#1\empty\else%
    \PYG@tok{#1}\expandafter\PYG@toks\fi}
\def\PYG@do#1{\PYG@bc{\PYG@tc{\PYG@ul{%
    \PYG@it{\PYG@bf{\PYG@ff{#1}}}}}}}
\def\PYG#1#2{\PYG@reset\PYG@toks#1+\relax+\PYG@do{#2}}


\@namedef{PYG@tok@bp}{\def\PYG@tc##1{\textcolor[HTML]{000000}{##1}}}
\@namedef{PYG@tok@c}{\def\PYG@tc##1{\textcolor[HTML]{336666}{##1}}\def\PYG@it##1{\textit{##1}}}
\@namedef{PYG@tok@c1}{\def\PYG@tc##1{\textcolor[HTML]{336666}{##1}}\def\PYG@it##1{\textit{##1}}}
\@namedef{PYG@tok@ch}{\def\PYG@tc##1{\textcolor[HTML]{336666}{##1}}\def\PYG@it##1{\textit{##1}}}
\@namedef{PYG@tok@cm}{\def\PYG@tc##1{\textcolor[HTML]{336666}{##1}}\def\PYG@it##1{\textit{##1}}}
\@namedef{PYG@tok@cp}{\def\PYG@tc##1{\textcolor[HTML]{336666}{##1}}\def\PYG@it##1{\textit{##1}}}
\@namedef{PYG@tok@cpf}{\def\PYG@tc##1{\textcolor[HTML]{336666}{##1}}\def\PYG@it##1{\textit{##1}}}
\@namedef{PYG@tok@cs}{\def\PYG@tc##1{\textcolor[HTML]{336666}{##1}}\def\PYG@it##1{\textit{##1}}}
\@namedef{PYG@tok@dl}{\def\PYG@tc##1{\textcolor[HTML]{000000}{##1}}}
\@namedef{PYG@tok@err}{\def\PYG@tc##1{\textcolor[HTML]{000000}{##1}}}
\@namedef{PYG@tok@fm}{\def\PYG@tc##1{\textcolor[HTML]{000000}{##1}}}
\@namedef{PYG@tok@gd}{\def\PYG@tc##1{\textcolor[HTML]{000000}{##1}}}
\@namedef{PYG@tok@ge}{\def\PYG@tc##1{\textcolor[HTML]{000000}{##1}}}
\@namedef{PYG@tok@ges}{\def\PYG@tc##1{\textcolor[HTML]{000000}{##1}}}
\@namedef{PYG@tok@gh}{\def\PYG@tc##1{\textcolor[HTML]{000000}{##1}}}
\@namedef{PYG@tok@gi}{\def\PYG@tc##1{\textcolor[HTML]{000000}{##1}}}
\@namedef{PYG@tok@go}{\def\PYG@tc##1{\textcolor[HTML]{000000}{##1}}}
\@namedef{PYG@tok@gp}{\def\PYG@tc##1{\textcolor[HTML]{000000}{##1}}}
\@namedef{PYG@tok@gr}{\def\PYG@tc##1{\textcolor[HTML]{000000}{##1}}}
\@namedef{PYG@tok@gs}{\def\PYG@tc##1{\textcolor[HTML]{000000}{##1}}}
\@namedef{PYG@tok@gt}{\def\PYG@tc##1{\textcolor[HTML]{000000}{##1}}}
\@namedef{PYG@tok@gu}{\def\PYG@tc##1{\textcolor[HTML]{000000}{##1}}}
\@namedef{PYG@tok@il}{\def\PYG@tc##1{\textcolor[HTML]{000000}{##1}}}
\@namedef{PYG@tok@k}{\def\PYG@tc##1{\textcolor[HTML]{006699}{##1}}\def\PYG@bf##1{\textbf{##1}}}
\@namedef{PYG@tok@kc}{\def\PYG@tc##1{\textcolor[HTML]{006699}{##1}}\def\PYG@bf##1{\textbf{##1}}}
\@namedef{PYG@tok@kd}{\def\PYG@tc##1{\textcolor[HTML]{006699}{##1}}\def\PYG@bf##1{\textbf{##1}}}
\@namedef{PYG@tok@kn}{\def\PYG@tc##1{\textcolor[HTML]{006699}{##1}}\def\PYG@bf##1{\textbf{##1}}}
\@namedef{PYG@tok@kp}{\def\PYG@tc##1{\textcolor[HTML]{006699}{##1}}\def\PYG@bf##1{\textbf{##1}}}
\@namedef{PYG@tok@kr}{\def\PYG@tc##1{\textcolor[HTML]{006699}{##1}}\def\PYG@bf##1{\textbf{##1}}}
\@namedef{PYG@tok@kt}{\def\PYG@tc##1{\textcolor[HTML]{007788}{##1}}\def\PYG@bf##1{\textbf{##1}}}
\@namedef{PYG@tok@kw}{\def\PYG@tc##1{\textcolor[HTML]{006699}{##1}}\def\PYG@bf##1{\textbf{##1}}}
\@namedef{PYG@tok@m}{\def\PYG@tc##1{\textcolor[HTML]{ff6600}{##1}}}
\@namedef{PYG@tok@mb}{\def\PYG@tc##1{\textcolor[HTML]{ff6600}{##1}}}
\@namedef{PYG@tok@mf}{\def\PYG@tc##1{\textcolor[HTML]{ff6600}{##1}}}
\@namedef{PYG@tok@mh}{\def\PYG@tc##1{\textcolor[HTML]{ff6600}{##1}}}
\@namedef{PYG@tok@mi}{\def\PYG@tc##1{\textcolor[HTML]{ff6600}{##1}}}
\@namedef{PYG@tok@mo}{\def\PYG@tc##1{\textcolor[HTML]{ff6600}{##1}}}
\@namedef{PYG@tok@na}{\def\PYG@tc##1{\textcolor[HTML]{007788}{##1}}\def\PYG@bf##1{\textbf{##1}}}
\@namedef{PYG@tok@nb}{\def\PYG@tc##1{\textcolor[HTML]{000088}{##1}}}
\@namedef{PYG@tok@nc}{\def\PYG@tc##1{\textcolor[HTML]{00aa88}{##1}}\def\PYG@bf##1{\textbf{##1}}}
\@namedef{PYG@tok@nd}{\def\PYG@tc##1{\textcolor[HTML]{00aa88}{##1}}\def\PYG@bf##1{\textbf{##1}}}
\@namedef{PYG@tok@ne}{\def\PYG@tc##1{\textcolor[HTML]{000000}{##1}}}
\@namedef{PYG@tok@nf}{\def\PYG@tc##1{\textcolor[HTML]{cc00ff}{##1}}}
\@namedef{PYG@tok@ni}{\def\PYG@tc##1{\textcolor[HTML]{000000}{##1}}}
\@namedef{PYG@tok@nl}{\def\PYG@tc##1{\textcolor[HTML]{000088}{##1}}}
\@namedef{PYG@tok@nn}{\def\PYG@tc##1{\textcolor[HTML]{000088}{##1}}}
\@namedef{PYG@tok@no}{\def\PYG@tc##1{\textcolor[HTML]{000088}{##1}}}
\@namedef{PYG@tok@nt}{\def\PYG@tc##1{\textcolor[HTML]{000000}{##1}}}
\@namedef{PYG@tok@nv}{\def\PYG@tc##1{\textcolor[HTML]{000088}{##1}}}
\@namedef{PYG@tok@o}{\def\PYG@tc##1{\textcolor[HTML]{555555}{##1}}}
\@namedef{PYG@tok@ow}{\def\PYG@tc##1{\textcolor[HTML]{555555}{##1}}}
\@namedef{PYG@tok@p}{\def\PYG@tc##1{\textcolor[HTML]{555555}{##1}}}
\@namedef{PYG@tok@s}{\def\PYG@tc##1{\textcolor[HTML]{cc3300}{##1}}}
\@namedef{PYG@tok@s1}{\def\PYG@tc##1{\textcolor[HTML]{cc3300}{##1}}}
\@namedef{PYG@tok@s2}{\def\PYG@tc##1{\textcolor[HTML]{cc3300}{##1}}}
\@namedef{PYG@tok@sa}{\def\PYG@tc##1{\textcolor[HTML]{cc3300}{##1}}}
\@namedef{PYG@tok@sb}{\def\PYG@tc##1{\textcolor[HTML]{cc3300}{##1}}}
\@namedef{PYG@tok@sc}{\def\PYG@tc##1{\textcolor[HTML]{cc3300}{##1}}}
\@namedef{PYG@tok@sd}{\def\PYG@tc##1{\textcolor[HTML]{cc3300}{##1}}}
\@namedef{PYG@tok@se}{\def\PYG@tc##1{\textcolor[HTML]{cc3300}{##1}}}
\@namedef{PYG@tok@sh}{\def\PYG@tc##1{\textcolor[HTML]{cc3300}{##1}}}
\@namedef{PYG@tok@si}{\def\PYG@tc##1{\textcolor[HTML]{000000}{##1}}}
\@namedef{PYG@tok@sr}{\def\PYG@tc##1{\textcolor[HTML]{cc3300}{##1}}}
\@namedef{PYG@tok@ss}{\def\PYG@tc##1{\textcolor[HTML]{cc3300}{##1}}}
\@namedef{PYG@tok@sx}{\def\PYG@tc##1{\textcolor[HTML]{cc3300}{##1}}}
\@namedef{PYG@tok@vc}{\def\PYG@tc##1{\textcolor[HTML]{000088}{##1}}}
\@namedef{PYG@tok@vg}{\def\PYG@tc##1{\textcolor[HTML]{000088}{##1}}}
\@namedef{PYG@tok@vi}{\def\PYG@tc##1{\textcolor[HTML]{000088}{##1}}}
\@namedef{PYG@tok@vm}{\def\PYG@tc##1{\textcolor[HTML]{000088}{##1}}}
\@namedef{PYG@tok@w}{\def\PYG@tc##1{\textcolor[HTML]{000000}{##1}}}

\def\PYGZbs{\char`\\}
\def\PYGZus{\char`\_}
\def\PYGZob{\char`\{}
\def\PYGZcb{\char`\}}
\def\PYGZca{\char`\^}
\def\PYGZam{\char`\&}
\def\PYGZlt{\char`\<}
\def\PYGZgt{\char`\>}
\def\PYGZsh{\char`\#}
\def\PYGZpc{\char`\%}
\def\PYGZdl{\char`\$}
\def\PYGZhy{\char`\-}
\def\PYGZsq{\char`\'}
\def\PYGZdq{\char`\"}
\def\PYGZti{\char`\~}
\def\PYGZat{@}
\def\PYGZlb{[}
\def\PYGZrb{]}
\makeatother
     % \PYG 系列宏（与 minted 的样式文件同一套）
\fvset{breaklines=true, breakanywhere=true, fontsize=\codefont,
       xleftmargin=17bp, tabsize=4, commandchars=\\\{\}}
\fi

% -----------------------------------------------------------------------------
%  标题样式（对齐原书：压缩无衬线、纯黑、右对齐章首）
%   章：CHAPTER N（16.8bp）+ 通栏细线 + 章名（25.2bp，右对齐）
%   节：18.9bp；小节：15.8bp —— 编号在标题行内，不悬挂（中文版保留编号）
%   章首版式不用 titlesec 的 [display] 形状：那里不允许 label 里出现 \par，
%   而「ChAPTER N / 细线 / 章名」三行结构必须在标题块内完成，故直接重定义
%   \@makeschapterhead（全书只用 \chapter*，编号由 \myChapter 手工设定）。
% -----------------------------------------------------------------------------
\usepackage{titlesec}
\newcommand{\mychaptergiant}{}
\newcommand{\mychapterset}[1]{\gdef\mychaptergiant{#1}}
\newcommand{\mychapterclear}{\gdef\mychaptergiant{}}
\newcommand{\myChapterTopSkip}{3bp}
\newcommand{\myChapterAfterSkip}{96bp}

\makeatletter
\def\@makeschapterhead#1{%
  \vspace*{\myChapterTopSkip}%
  \thispagestyle{plain}%
  \begingroup
    \raggedleft\sffamily
    \ifx\mychaptergiant\empty\else
      {\fontsize{16.8bp}{20bp}\selectfont\bfseries CHAPTER~\mychaptergiant\par}%
      \vspace{1bp}%
    \fi
    {\color{ink}\rule{\textwidth}{0.5bp}\par}%
    \vspace{-10bp}%
    {\fontsize{25.2bp}{28bp}\selectfont\bfseries #1\par}%
  \endgroup
  \vspace{\myChapterAfterSkip}%
  \@afterheading}
\makeatother

\titleformat{\section}
  {\fontsize{18.9bp}{22bp}\selectfont\sffamily\bfseries\color{ink}}
  {\thesection\hspace{0.5em}}{0bp}{}
\titlespacing*{\section}{0bp}{22bp plus 4bp minus 2bp}{7bp}

\titleformat{\subsection}
  {\fontsize{15.8bp}{19bp}\selectfont\sffamily\bfseries\color{ink}}
  {\thesubsection\hspace{0.5em}}{0bp}{}
\titlespacing*{\subsection}{0bp}{16bp plus 3bp minus 2bp}{5bp}

\setcounter{secnumdepth}{2}     % 章 / 节 / 小节都编号（对齐中文原版）
\setcounter{tocdepth}{2}        % 目录与书签只到「小节」

% -----------------------------------------------------------------------------
%  插图：通栏细边框 + 下方题注（题注左缘与正文齐平）
%  用法：\myFigure{images/ch03/fig-3-1.jpeg}{图 3-1：…}
% -----------------------------------------------------------------------------
%  ⚠ 图与题注必须打进同一个 \vbox：若用 center 环境 + \nopagebreak，
%    \end{center} 之后的 \vspace 本身就是合法断点，图排到页底时题注会
%    单独跑到下一页（第 3 章图 3-3 实测）。整块当成一个垂直盒子就拆不开。
\newcommand{\myFigure}[2]{%
  \par\addvspace{9bp}%
  \noindent
  \vbox{%
    \hbox to \textwidth{\hfil
      {\setlength{\fboxrule}{0.4bp}\setlength{\fboxsep}{0pt}%
       \fcolorbox{figframe}{white}{%
         \includegraphics[max width=\textwidth, max height=0.52\textheight,
                          keepaspectratio]{#1}}}%
      \hfil}%
    \ifx\relax#2\relax\else
      \vskip5bp
      \hbox to \textwidth{\vtop{\hsize=\textwidth
        \fontsize{9.8bp}{12.5bp}\selectfont\bfseries #2\par}\hfil}%
    \fi}%
  \par\addvspace{9bp}}

% 表/图题注（原文单独成段加粗）：与图注同一套字号，段前段后收紧
\newcommand{\myCaption}[1]{%
  \par\addvspace{8bp}\noindent
  {\fontsize{9.8bp}{12.5bp}\selectfont\bfseries #1\par}%
  \addvspace{3bp}}

% -----------------------------------------------------------------------------
%  表格：无竖线、表头黑底白字（对齐原书）
%  ⚠ 不能把 tabularx 包进自定义环境：tabularx 要**向前扫描**自己的 \end{tabularx}
%    才能确定列宽，包一层会让它读到 \end{myTable}，报
%    "File ended while scanning use of \TX@get@body"。
%    所以只在两侧放「开关宏」，tabularx 本体留在正文里。
%    行数多的表（>=14 行）由 convert.py 改用 xltabular，以便跨页并重复表头：
%    tabularx 是不可分割的段落，比一页高的表会越过页脚、多出的行静默丢失。
% -----------------------------------------------------------------------------
\newcommand{\myTableBegin}{%
  \par\addvspace{2bp}\noindent\begingroup
  \fontsize{9bp}{11.5bp}\selectfont\sffamily
  \renewcommand{\arraystretch}{1.25}%
  \setlength{\arrayrulewidth}{0.3bp}%
  % longtable 默认在表格前后各加 \bigskipamount，且靠 \fill 居中；
  % 这里归零，让表与正文左缘严格对齐、段距只由 \myTableBegin/\myTableEnd 控制。
  \setlength{\LTpre}{0bp}\setlength{\LTpost}{0bp}%
  \setlength{\LTleft}{0bp}\setlength{\LTright}{0bp}}
\newcommand{\myTableEnd}{\endgroup\par\addvspace{9bp}}

% -----------------------------------------------------------------------------
%  提示框（NOTE / TIP / WARNING）：页边距里的图标 + 缩进正文，无边框无底色
%  原书实测：图标位于版心左缘外 17bp、宽 38bp；正文整块缩进 65bp，字号 9.6bp
% -----------------------------------------------------------------------------
\newcommand{\myNoteIcon}{images/icon-note.png}
\newenvironment{myNote}
  {\par\addvspace{8bp}\begingroup
   \setlength{\leftskip}{65bp}\setlength{\parindent}{0bp}%
   \noindent\hspace*{-65bp}%
   \raisebox{-4bp}{\includegraphics[width=38bp]{\myNoteIcon}}%
   \hspace*{27bp}%
   \fontsize{9.6bp}{12.6bp}\selectfont\ignorespaces}
  {\par\endgroup\addvspace{8bp}}

% 短小的旁注（原文中的缩进小段）
\newenvironment{mySidenote}
  {\par\addvspace{5bp}\begingroup
   \setlength{\leftskip}{2.2em}\setlength{\parindent}{0bp}%
   \fontsize{9.6bp}{12.6bp}\selectfont\ignorespaces}
  {\par\endgroup\addvspace{5bp}}

% 章首/节首引文（epigraph）
%   正文斜体 9.5bp、左侧缩进 18bp；出处另起一行、同样缩进、不斜体
%   注意：环境的结束码里拿不到 #1，出处必须先存进 \my@epiattr。
\newcommand{\myepiattr}{}
\newenvironment{myEpigraph}[1]
  {\par\addvspace{10bp}\begingroup
   \gdef\myepiattr{#1}%
   \setlength{\leftskip}{18bp}\setlength{\parindent}{0bp}%
   \itshape\fontsize{9.5bp}{12.4bp}\selectfont\noindent\ignorespaces}
  {\par\nopagebreak\vspace{2bp}%
   \upshape\fontsize{9bp}{11.5bp}\selectfont\noindent
   \ifx\myepiattr\empty\else\myepiattr\fi
   \par\endgroup\addvspace{10bp}}

% -----------------------------------------------------------------------------
%  目录 / 书签 / 超链接
% -----------------------------------------------------------------------------
\usepackage{url}
\usepackage[depth=2]{bookmark}
\usepackage{hyperref}
\hypersetup{
  colorlinks=true,
  linkcolor=linkred,
  citecolor=linkred,
  urlcolor=linkred,
  filecolor=linkred,
  breaklinks=true,
  pdfstartview=Fit,
  pdftitle={Rust 程序设计（第 2 版）},
  pdfauthor={Jim Blandy, Jason Orendorff, Leonora F. S. Tindall 著；汪志成 译},
  pdfsubject={Programming Rust: Fast, Safe Systems Development, Second Edition 中文版},
  pdfkeywords={Rust, 系统编程, 所有权, 生命周期, 并发, 异步, 宏, FFI},
}
\renewcommand{\contentsname}{目录}

% 章名上的脚注标记。直接用 \footnotemark 有两个问题：被 hyperref 包成链接、
% 染成暗红；以及 \@makefnmark 的 7bp 是按正文 10.5bp 定的，放进 25.2bp 的
% 章名里几乎看不见。这里局部改回正文字色，并给标记 11bp。
\makeatletter
\newcommand{\mychapterfnmark}{{\hypersetup{linkcolor=ink}%
  \def\@makefnmark{\hbox{\textsuperscript{\fontsize{11bp}{11bp}\selectfont\@thefnmark}}}%
  \footnotemark}}
\makeatother

% 目录页：标题后的空白比章首页小得多（原书目录紧跟标题），
% 且目录条目不跟着链接染成暗红——书里的目录是纯黑。
\makeatletter
\let\myORItableofcontents\tableofcontents
\renewcommand{\tableofcontents}{%
  \begingroup
    \hypersetup{linkcolor=ink}%
    \renewcommand{\myChapterAfterSkip}{44bp}%
    \myORItableofcontents
  \endgroup}
\makeatother

% 目录条目：章条目补引导点，页码盒加宽到 1.9em（避免三位页码 Overfull）
\makeatletter
\def\@pnumwidth{1.9em}
\renewcommand*\l@chapter[2]{%
  \ifnum \c@tocdepth >\m@ne
    \addpenalty{-\@highpenalty}%
    \vskip 0.9em \@plus\p@
    \begingroup
      \parindent \z@ \rightskip \@pnumwidth
      \parfillskip -\@pnumwidth
      \leavevmode \bfseries
      #1\nobreak
      \leaders\hbox{$\m@th\mkern\@dotsep mu\hbox{.}\mkern\@dotsep mu$}\hfill
      \nobreak\hb@xt@\@pnumwidth{\hss #2}\par
      \penalty\@highpenalty
    \endgroup
  \fi}
% 页脚里的节名：带节号（"2.6　并发"），去掉 book.cls 自带的句点
\newcommand{\my@sectionmark}[1]{\markright{\normalfont
  \ifnum\c@secnumdepth>\z@\thesection\quad\fi #1\quad|\quad}}
\let\sectionmark\my@sectionmark
\AtBeginDocument{\let\sectionmark\my@sectionmark}
\makeatother

% -----------------------------------------------------------------------------
%  长串断行：URL 与行内长标识符允许折行（同 CMake 模板的做法）
%   在 / . - _ ? = 之后放自然断点 \penalty 0，其余字符间放 \penalty 1000
%   应急断点；控制序列后不注入（否则 \emph 等宏会把 \penalty 当参数吃掉）。
% -----------------------------------------------------------------------------
\ExplSyntaxOn
\let\myORIhref\href
\let\myORItt\texttt
\cs_new:Npn \my_spcmark {}
\tl_new:N \l_myin_tl
\tl_new:N \l_myout_tl
\str_new:N \l_my_head_str
\str_new:N \l_my_first_str
\cs_new_protected:Npn \my_breakinject:n #1
  {
    \tl_set:Nn \l_myin_tl { #1 }
    \tl_replace_all:Nnn \l_myin_tl { ~ } { \my_spcmark }
    \tl_clear:N \l_myout_tl
    \tl_map_inline:Nn \l_myin_tl
      {
        \tl_put_right:Nn \l_myout_tl { ##1 }
        \str_set:Nn \l_my_head_str { ##1 }
        \str_set:Nx \l_my_first_str { \str_item:Nn \l_my_head_str { 1 } }
        \str_if_eq:VVTF \l_my_first_str \c_backslash_str
          { \tl_if_eq:nnT { ##1 } { \_ }
              { \tl_put_right:Nn \l_myout_tl { \penalty 0 \scan_stop: } } }
          { \tl_put_right:Nn \l_myout_tl { \penalty 1000 \scan_stop: } }
      }
    \tl_replace_all:Nnn \l_myout_tl { / \penalty 1000 \scan_stop: } { / \penalty 0 \scan_stop: }
    \tl_replace_all:Nnn \l_myout_tl { . \penalty 1000 \scan_stop: } { . \penalty 0 \scan_stop: }
    \tl_replace_all:Nnn \l_myout_tl { - \penalty 1000 \scan_stop: } { - \penalty 0 \scan_stop: }
    \tl_replace_all:Nnn \l_myout_tl { ? \penalty 1000 \scan_stop: } { ? \penalty 0 \scan_stop: }
    \tl_replace_all:Nnn \l_myout_tl { = \penalty 1000 \scan_stop: } { = \penalty 0 \scan_stop: }
    \tl_replace_all:Nnn \l_myout_tl { \_ \penalty 1000 \scan_stop: } { \_ \penalty 0 \scan_stop: }
    \tl_replace_all:Nnn \l_myout_tl { \my_spcmark } { ~ }
  }
\RenewDocumentCommand { \href } { m m }
  { \my_breakinject:n { #2 } \myORIhref { #1 } { \l_myout_tl } }
\RenewDocumentCommand { \texttt } { m }
  { \my_breakinject:n { #1 } \myORItt { \l_myout_tl } }
\ExplSyntaxOff

% 裸 URL（原文链接文字就是网址）：链接目标 = 显示文字
\newcommand{\texthref}[1]{\href{#1}{#1}}

% -----------------------------------------------------------------------------
%  页脚（对齐原书）：双页 = 「页码 | 章名」靠左；单页 = 「节名 | 页码」靠右；
%  上方一条 0.3bp 细线；章首页只印页码，位置在外侧。
% -----------------------------------------------------------------------------
\usepackage{fancyhdr}
\pagestyle{fancy}
\fancyhf{}
\newcommand{\myfootfont}{\sffamily\fontsize{9bp}{11bp}\selectfont}
% 分隔符（竖线）跟着书眉内容走，而不是写死在页脚里：
%   书眉为空（封面/扉页/版权页等 \markboth{}{} 的页）时自然什么都不印，
%   否则会留下一个孤零零的「| 3」。
%   也不能用 etoolbox 的 \ifdefempty 去判 \@leftmark —— LaTeX 在标记里
%   存了大小写包装，永远判不出"空"。
\fancyfoot[LE]{\myfootfont\bfseries\thepage\leftmark}
\fancyfoot[RO]{\myfootfont\rightmark\bfseries\thepage}
\renewcommand{\headrulewidth}{0bp}
\renewcommand{\footrulewidth}{0.3bp}
\renewcommand{\footrule}{%
  \vskip-\footruleskip\vskip-\footrulewidth
  {\color{rule}\hrule height\footrulewidth width\headwidth}%
  \vskip\footruleskip}
\fancypagestyle{plain}{%
  \fancyhf{}%
  \fancyfoot[LE]{\myfootfont\bfseries\thepage}%
  \fancyfoot[RO]{\myfootfont\bfseries\thepage}%
  \renewcommand{\headrulewidth}{0bp}%
  \renewcommand{\footrulewidth}{0.3bp}}
% 书眉内容自带分隔符：左页「章名」、右页「章名/节名」，竖线由 \chapmark 与
% \my@sectionmark 一并写进去，页脚只负责拼接（见上一条说明）。
\newcommand{\chapmark}[1]{%
  \markboth{\normalfont\quad|\quad #1}{\normalfont #1\quad|\quad}}

% 空白页（openright 自动插入的偶数页）不印页脚：
%   book.cls 的 \cleardoublepage 不会给空白页设 \pagestyle{empty}，fancy 页脚
%   会留下「12 |」这种只有页码和竖线的残缺行。这里补一个只对空白页生效的 empty。
\let\myorigcleardoublepage\cleardoublepage
\renewcommand{\cleardoublepage}{%
  \clearpage{\pagestyle{empty}\myorigcleardoublepage}}

% -----------------------------------------------------------------------------
%  Unicode 符号兜底（当前字体缺字形时映射到等价写法）
% -----------------------------------------------------------------------------
\usepackage{newunicodechar}
\newunicodechar{→}{\(\to\)}
\newunicodechar{←}{\(\leftarrow\)}
\newunicodechar{⇒}{\(\Rightarrow\)}
\newunicodechar{≤}{\(\leq\)}
\newunicodechar{≥}{\(\geq\)}
\newunicodechar{≠}{\(\neq\)}
\newunicodechar{−}{\(-\)}
\newunicodechar{×}{\(\times\)}
\newunicodechar{·}{\(\cdot\)}
\newunicodechar{∞}{\(\infty\)}
\newunicodechar{⊂}{\(\subset\)}
\newunicodechar{⊆}{\(\subseteq\)}
\newunicodechar{∈}{\(\in\)}
\newunicodechar{∠}{\(\angle\)}
\newunicodechar{⊙}{\(\odot\)}
\newunicodechar{≅}{\(\cong\)}
\newunicodechar{π}{\(\pi\)}
\newunicodechar{μ}{\(\mu\)}
\newunicodechar{β}{\(\beta\)}
\newunicodechar{✓}{\(\checkmark\)}
\newunicodechar{✗}{\(\times\)}
\newunicodechar{➊}{\textcircled{\scriptsize 1}}
\newunicodechar{➋}{\textcircled{\scriptsize 2}}
\newunicodechar{➌}{\textcircled{\scriptsize 3}}
\newunicodechar{■}{\rule[0.2em]{0.62em}{0.62em}}
\newunicodechar{▪}{\rule[0.2em]{0.4em}{0.4em}}
\newunicodechar{☐}{\(\square\)}
\newunicodechar{️}{}

% 非拉丁字符的字体回退（Cambria / Consolas 都没有这些字形）
%   · 第 17 章讲 String::chars 时的举例里有希伯来文 "ערב טוב"
%   · 第 3 章讲字符类型时的举例里有卡纳达文 "ಠ_ಠ"
%   · 第 18 章讲 Path 转字符串时会提到替换字符 U+FFFD
% 代码块里的回退由 tools/convert.py 在生成 \PYG 序列时包 {\hebrewfont …} 完成
% （Verbatim 的 commandchars 允许花括号与反斜杠，故在代码块内同样生效）。
\newfontfamily\hebrewfont{Arial}
\newfontfamily\kannadafont{Nirmala UI}
\newfontfamily\symbolfallback{Segoe UI Symbol}
% U+FFFD（替换字符）没有一款 Windows 字体真正带这个码位，原书里它是一个空心方框，
% 这里直接画一个，既忠实又不留豆腐块。
\newunicodechar{�}{{\setlength{\fboxsep}{0pt}%
  \framebox[0.62em]{\rule{0pt}{0.5em}}}}
\newunicodechar{ಠ}{{\kannadafont ಠ}}

% -----------------------------------------------------------------------------
%  章节命令（对应模板的 \myChapter / \myFront）
% -----------------------------------------------------------------------------
\ExplSyntaxOn
% 编号正文章节：\myChapter{章号}{标题}{文件}
\tl_new:N \l__my_toctitle_tl
\NewDocumentCommand { \myChapter } { mmm }
{
  \setcounter{chapter}{#1}
  \setcounter{section}{0}
  % 脚注按章从 1 重新编（原书就是每章独立编号）。不重置的话，
  % 第 4 章章名上那条译者注会显示成 [14]，因为前面三章累计了 13 条。
  \setcounter{footnote}{0}
  \mychapterset{#1}%
  \tl_set:Nn \l__my_toctitle_tl { #2 }
  % 章名上可能挂 \footnotemark（原书章名带脚注），它不该进目录与页眉
  \tl_remove_all:Nn \l__my_toctitle_tl { \mychapterfnmark }
  \tl_replace_all:Nnn \l__my_toctitle_tl { \\ } { \space }
  \chapter*{#2}
  \mychapterclear
  \chapmark{第~#1~章\quad \l__my_toctitle_tl}
  \addcontentsline{toc}{chapter}{第~#1~章\quad \l__my_toctitle_tl}
  \subfile{#3}
}
% 前言类：进目录但无编号（章首无「CHAPTER N」行，标题上方无细线，
%   标题后不留章首页那样的大空白）
\NewDocumentCommand { \myFront } { mm }
{
  \mychapterclear
  \setcounter{footnote}{0}
  \group_begin:
    \cs_set:Npn \myChapterAfterSkip { 34bp }
    \cs_set:Npn \myChapterTopSkip { 24bp }
    \chapter*{#1}
  \group_end:
  \addcontentsline{toc}{chapter}{#1}
  \chapmark{#1}
  % 前言类的小节不编号，否则目录里会冒出「0.1 致谢」「0.2 读者对象」这种
  % 挂在第 0 章下的编号。\setcounter 是全局赋值，正文前再恢复成 2。
  \setcounter{secnumdepth}{-1}
  \subfile{#2}
  \setcounter{secnumdepth}{2}
}
% 版权页：不进目录
\NewDocumentCommand { \myFrontNoToc } { m }
{
  \clearpage
  \markboth{}{}
  \subfile{#1}
}
\ExplSyntaxOff

% 版权页用的小号排版
\newenvironment{myCopyright}
  {\par\begingroup\fontsize{9bp}{13bp}\selectfont
   \setlength{\parindent}{0bp}\setlength{\parskip}{4bp}}
  {\par\endgroup}
